\section{Nomination persistante des orbites}
\label{sec:reevaluation-nomination-persistante-orbites}

\maxime{idée plan~:\\
	Combinaison d'orbites d'origine et de suivi dans un DAG pour une représentation
	complète d'un évènement à chaque application de règles. Formalisation de l'usage
	de chemins pour préciser une origine lorsque celle-ci n'est pas directement liée
	à un nœud d'accroche dans une règle. }

Pour la deuxième étape de notre nomination, nous représentons l'histoire de
chacune des entités topologiques, désignées par des noms persistants de brins,
dans la spécification paramétrique. %
Cette histoire, nous la représentons à l'aide d'un graphe orienté acyclique
(DAG en anglais) qui retrace les changements subis par les ancêtres d'une
entité. %

\maxime{Ici, il y a plusieurs concepts qui interviennent~:\\
	\begin{itemize}
		\item les entrées pour la construction du DAG sont~: un nom persistant de
		      brin, le type d'orbite de l'entité topologique, une spécification
		      paramétrique (pour associer numéro d'opération et opération)
		\item les notions de suivi (flèche noire) et d'origine (flèche rouge)
		      d'une orbite (au travers d'une opération)
	\end{itemize}
}

\maxime{On ressort la construction fil rouge, et on illustre l'histoire de la
	face coloriée à l'issue de la dernière opération du processus de modélisation}

\maxime{1. On commence par désigner le nom persistant de brin à partir duquel on
	va représenter une entité topologique} \maxime{2. On montre le type d'orbite
	associé au hook de la règle $\to$ désignation d'une entité topologique}

\maxime{3. Tracé de l'histoire à partir de l'entité désignée à l'étape
	précédente (étape finale du processus de modélisation) jusqu'aux ancêtres les
	plus anciens (potentiellement la première étape du processus de modélisation)}

\maxime{4. étapes successives de restitution de l'histoire d'une entité
	topologique à partir d'illustrations. Faire attention de bien présenter qu'il
	y a à la fois le suivi et les origines. Faire attention à bien expliquer que
	la notion d'origine a une signification dans la restitution et la lecture de
	l'histoire}

\maxime{5. On ajoute une notion qui est celle de chemin~: parfois, le lien entre
	une orbite et son origine ne peut être fait de manière directe parce que
	l'origine a été supprimée dans l'opération. %
	Dans ce cas, on peut annoter une orbite avec un chemin qui fait le lien entre
	l'origine et les ancêtres directs d'une orbite.}

À terme, un DAG d'évaluation représente l'histoire d'une unique entité topologique. %
Les apports de cette approche sont~:
\begin{itemize}
	\item la détection statique des évènements (formalisés dans le chapitre
	      précédent) permet de reconstituer l'histoire d'une entité
	\item pas de besoin de réappliquer les règles, c'est-à-dire que pas besoin
	      de reproduire la représentation explicite
	\item uniquement la spécification paramétrique
	\item uniquement pour les entités calculée à partir du type d'orbite d'un
	      hook et du nom persistant de brin associé dans la spécification paramétrique.
\end{itemize}


%%% Local Variables:
%%% mode: LaTeX
%%% TeX-master: "../../../main"
%%% End:
